% %--------------------------
% %	CONFIGURATION
% %--------------------------
\documentclass[11pt,a4paper]{moderncv}
\moderncvstyle{classic}
\moderncvcolor{black}
\definecolor{firstnamecolor}{rgb}{0,0,0}

% Basic packages
\usepackage[utf8]{inputenc}
\usepackage[T1]{fontenc}
\usepackage[hscale=0.7, top=2cm, bottom=2cm]{geometry}
\usepackage[dvipsnames]{xcolor}
\definecolor{DarkBlue}{RGB}{0, 0, 150}
\usepackage[colorlinks=true, urlcolor=DarkBlue, linkcolor=DarkBlue]{hyperref}

\usepackage{microtype}

% Configure font settings
\renewcommand{\rmdefault}{ppl}
\renewcommand{\sfdefault}{phv}
\renewcommand*{\namefont}{\fontsize{26}{18}\mdseries}

% Configure microtype
\microtypesetup{expansion=false}

% Load sensitive information
\input{../../other_material/personal_info_dk.tex}

% %--------------------------
\begin{document}
\makecvtitle
Hello Rodrigo Thomas,

\vspace*{2em}
I am writing to apply for the Quantum Photonics Engineer role. I am not the standard candidate for it as my recent professional work is data science and ML engineering. I do however have direct optics-lab experience from an exchange at Zhejiang University (ZJU), where I worked with a local research group on fabricating microring resonators and tuning their resonance frequency with controlled laser pulses to selectively remove cladding material. I am applying because I am eager to work on the cutting edge of research and Quantum Foundry Copenhagen's work on quantum processor manufacturing aligns well with my background in physics.

\hspace*{2em}
I have a background in Engineering Physics from the Royal Institute of Technology (KTH), followed by a master's specialisation in Statistical Learning and Data Analytics. I graduated with a 5.0/5.0 GPA, with coursework including Quantum Physics, Solid State Physics, and Electromagnetic Theory as well as courses in machine learning. My work at ZJU gave me practical exposure to integrated photonic circuits, optical alignment, laser-based device modification, lab instrumentation, and the importance of repeatable experimental protocols.

\hspace*{2em}
Professionally, I have worked on reliable applied ML and data systems. At Signum Life Science, I am building a Databricks forecasting platform for pharmaceutical price and demand. At Valcon, I architected an end-to-end production pipeline for a real-time neural network and led projects across pharma, energy, IoT, and manufacturing. Earlier, at RaySearch Laboratories, I worked on automated radiotherapy planning using autoencoders, sparse Gaussian processes, and optimisation on segmented 3D CT scans, in close collaboration with medical physicists and clinicians.

\hspace*{2em}
Though I have much to learn about quantum computing, and quantum foundry, I am confident I can keep up with the steep learning curve. My strong physics background will be an immense help in this but I also bring many relevant skills from within the other disciplines you mention on your website, such as mechanical engineering, electrical engineering, quantum physics, and data science. I am particularly interested in the characterization side of the role: developing repeatable measurement protocols, automating data acquisition in Python, and extracting reliable device-level performance metrics.

\hspace*{2em}
Outside work, I spend much of my time designing and building physical objects through woodworking and 3D printing. This reflects the same practical, curious and creative engineering mindset I would bring to experimental photonics work.

\vspace*{2em}
Thank you for your consideration. I hope to have the opportunity to discuss my application further. Regards,

\vspace*{1em}
Ivar Cashin Eriksson.

\end{document}
